\section*{Methods}

\subsection*{Cohort Definition and Important Feature Discovery}

Individuals were stratified into five non-overlapping analytical cohorts on the basis of their polygenic risk score (PRS) profile across established type~2 diabetes (T2D) genetic architectures.
Within each cohort, exclusive high-risk cases were defined as individuals scoring in the top 20\,\% of the PRS distribution for that cohort but \emph{not} exceeding the top-20\,\% threshold in any other cohort.
Within-cohort controls comprised individuals scoring in the bottom 20\,\% of the same PRS distribution with a confirmed non-T2D phenotype.

Individuals were assigned to three training strata, each subject to directional quality filters applied prior to enrichment:
\begin{itemize}
  \item \textbf{HighCases}: top-20\,\% T2D cases; only features with OR\,$>$\,1 (risk-elevating direction, $\beta > 0$) were retained.
  \item \textbf{LowControls}: bottom-20\,\% non-cases; only features with OR\,$<$\,1 (protective direction, $\beta < 0$) were retained. Risk-direction features ($\beta \geq 0$) were excluded to prevent contamination of the protective signal by poorly fitting model coefficients.
  \item \textbf{ExclusiveHighCases}: top-20\,\% T2D cases in exactly one cohort (i.e.\ \texttt{n\_exclusive\_cohorts}\,$= 1$) with OR\,$>$\,1; OR was derived as $\exp(\hat\beta)$ when not directly available from model output.
\end{itemize}

Important features were identified separately for two data modalities:

\paragraph{Genomic features.}
Cohort-exclusive genomic features were derived from pre-computed SHAP (SHapley Additive exPlanations) values obtained from gradient-boosted ensemble classifiers trained to discriminate high-risk cases from controls.
A feature was retained if it received a non-zero SHAP contribution in the high-risk case stratum (\texttt{HighCases} or \texttt{LowControls} training partitions, or the exclusive high-risk stratum \texttt{ExclusiveHighCases}).
The magnitude of each genomic feature's effect was summarised by its odds ratio (OR), derived as $\exp(\hat\beta)$ from the model logistic coefficient.
Features were classified as \texttt{genomic\_shap} if the feature identifier matched a dbSNP rs-identifier pattern or an HLA allele notation (e.g.\ \texttt{HLA-DRB1*401}, \texttt{HLA-B*4002}), and as \texttt{clinical} otherwise.

\paragraph{Clinical features.}
The importance of continuous clinical variables was assessed using two complementary Mann--Whitney~U frameworks.

\textit{Enrichment analysis} (\texttt{enrich\_exclusive\_cohort\_features.py}): clinical features were tested by comparing \texttt{ExclusiveHighCases} against within-cohort bottom-20\,\% controls using a two-sided Mann--Whitney~U test.
Both unbalanced (all controls) and balanced (bootstrap-resampled controls matched to $n_\text{cases}$) designs were applied; the median $p$-value across bootstrap replicates was used as the final significance statistic in the balanced mode.
Effect size was quantified as the rank-biserial correlation coefficient $r$.
Features were retained when Benjamini--Hochberg-corrected $p < 0.05$ and $|r| \geq 0.1$, and further filtered for cross-cohort specificity (see below).

\textit{Cross-cohort clinical comparison} (\texttt{compare\_clinical\_features\_across\_cohorts.py}): to characterise cohort-level clinical phenotype differences more broadly, the same Mann--Whitney~U framework was applied across three risk tiers:
\begin{itemize}
  \item \textbf{HighRisk}: each top-20\,\% cohort versus all remaining top-20\,\% cohorts combined (each-vs-others design);
  \item \textbf{LowRisk}: each bottom-20\,\% cohort versus all remaining bottom-20\,\% cohorts combined;
  \item \textbf{ExclusiveHighRisk}: exclusive top-20\,\% cases (one cohort only) versus within-cohort bottom-20\,\% controls, matching the case-vs-control design of the enrichment analysis.
\end{itemize}
FDR correction was applied independently within each tier.
A priority output was generated in which ExclusiveHighRisk results supersede HighRisk results for any feature--cohort pair that is significant in both, preserving the higher-specificity comparison where available.

\paragraph{Cross-cohort specificity tiering.}
To identify features whose importance was cohort-specific rather than globally elevated, a three-tier specificity filter was applied.
Tier~1 (strict): a feature was significant in the target cohort and not significant in any other cohort.
Tier~2 (majority): a feature was significant in the target cohort and non-significant in at least 60\,\% of the remaining cohorts.
Tier~3 (differential): a feature was significant in the target cohort, which exhibited the highest $|\text{effect}|$ across all cohorts, with a minimum advantage of $\Delta r \geq 0.1$ over the mean of other cohorts.
Each feature was assigned exclusively to the cohort with the highest $|r|$ or OR among all tier-qualifying cohorts.
Across the five analytical cohorts, between 4 and 333 features per cohort passed these criteria, spanning genomic (SHAP-derived, \texttt{genomic\_shap}) and clinical (\texttt{clinical}) data modalities.

\subsection*{Functional Cluster Enrichment}

\paragraph{Published functional cluster definitions.}
Functional clusters of T2D-associated loci and clinical biomarkers were sourced from four published resources:
\citet{Udler2018} defined five clusters of T2D loci (beta-cell function, proinsulin processing, obesity, lipodystrophy, and liver/lipid metabolism) using Bayesian non-negative matrix factorisation of GWAS summary statistics;
\citet{Kim2023} extended this framework to 16 clusters using multi-trait GWAS and co-localisation analyses, adding clusters for hyper-insulin secretion, lipoprotein~A, alkaline phosphatase, and SHBG;
\citet{Suzuki2024} applied a genome-wide signal aggregation approach identifying 380 distinct T2D signals grouped into eight functional clusters including a novel residual glycaemic cluster comprising the largest number of signals ($n = 389$), body fat, and metabolic syndrome;
\citet{Smith2024} performed a multi-ancestry meta-analysis extending the functional cluster architecture with additional clusters for lipodystrophy subtype~2, cholesterol, and bilirubin.
Clinical phenotype associations for each cluster were additionally informed by \citet{Mahajan2018} and \citet{Popovic2023}, who characterised the relationship between individual loci and quantitative clinical traits (HbA\textsubscript{1c}, fasting glucose, BMI, HOMA-B, HOMA-IR, triglycerides, HDL-cholesterol, and others).

\paragraph{Feature-to-cluster assignment.}
For each important feature identified in each cohort, cluster membership was determined by cross-referencing binary cluster indicator columns derived from the above publications.
Separate indicator matrices were constructed for loci maps and clinical phenotype maps from Udler~\textit{et al.}, Kim~\textit{et al.}, and Suzuki~\textit{et al.}
Cluster column names were normalised by converting to lower case, replacing hyphens, slashes, and underscores with spaces, removing the conjunction ``and'', inserting spaces before digit characters (e.g.\ \textit{beta cell1} $\to$ \textit{beta cell 1}), and collapsing consecutive whitespace.
Columns were matched to a curated \texttt{cluster\_subtype\_map.csv} by exact string match followed by a partial-substring fallback.
Genomic features (rs-identifiers) were matched to cluster membership via their nearest annotated gene symbol, prioritising the biological candidate gene (typically the second of two annotated flanking genes) over index-SNP neighbours such as \texttt{LOC}-prefix pseudogenes, \texttt{LINC}-prefix non-coding RNA genes, and antisense transcripts.
Clinical features were matched to cluster membership via a Mahajan~\textit{et al.}\ clinical phenotype bridge table that maps clinical variable names to the loci with which they are associated.
HLA-region alleles were classified as \texttt{genomic\_shap} features (not clinical features) and matched by gene-prefix; the classifier accepts both hyphen-delimited (\texttt{HLA-DQB1*504}) and underscore-delimited (\texttt{HLA\_DQB1\_504}) notations, as well as asterisk-containing allele designations (e.g.\ \texttt{HLA-drb1*401}), via the regular expression \texttt{\^{}HLA[-\_*]} applied case-insensitively.

\paragraph{Enrichment normalisation.}
For each cohort $c$ and functional cluster $k$, cluster enrichment was quantified as the number of cohort-exclusive important features mapping to that cluster:
\begin{equation}
  E_{c,k} = \left|\{f \in \mathcal{F}_c : f \mapsto k\}\right|,
\end{equation}
where $\mathcal{F}_c$ is the set of important features in cohort $c$.
Two normalisation schemes were applied to render counts comparable across cohorts of different sizes and clusters of different sizes:

\begin{enumerate}
  \item \textbf{Normalisation by cohort size}: the fraction of a cohort's total important features that map to cluster $k$,
  \begin{equation}
    \tilde{E}_{c,k}^{\text{cohort}} = \frac{E_{c,k}}{|\mathcal{F}_c|};
  \end{equation}

  \item \textbf{Normalisation by cluster size}: the fraction of a cluster's known constituent loci that are represented among cohort $c$'s important features,
  \begin{equation}
    \tilde{E}_{c,k}^{\text{cluster}} = \frac{E_{c,k}}{G_{c,k}},
  \end{equation}
  where $G_{c,k}$ is the total number of cluster-$k$ genes observed in cohort $c$.
\end{enumerate}

Bar heights in the circos plots reflect $|\log_2(\text{OR})|$ for genomic loci or $|r|$ for clinical features, scaled to the within-cohort maximum.
Bars extend outward when the association is risk-elevating ($\text{OR} > 1$ or $r > 0$) and inward when protective ($\text{OR} < 1$ or $r < 0$); inward bars are rendered with diagonal hatching (\texttt{//}) to distinguish them visually from outward bars of the same colour.
Genomic features are coloured red and clinical features yellow-green, regardless of direction.
A risk-only variant of each plot omits inward bars and protective (\texttt{LowControls}) loci entirely, showing only outward genomic and positive-effect clinical bars.

\paragraph{Directional concordance filtering.}
Each feature--cluster pair was assigned a concordance score based on the agreement between the observed direction of association and the published phenotypic profile of the corresponding functional cluster.
For clinical features, the observed direction was defined by the sign of the rank-biserial correlation $r$: a positive value indicates elevation of the trait in high-risk cases relative to controls; a negative value indicates reduction.
Published cluster phenotype profiles (derived from \citet{Udler2018}, \citet{Suzuki2024}, \citet{Kim2023}, and \citet{Yaghootkar2014}) were used to construct an expected-direction map specifying whether each measurable clinical trait should be elevated (+1) or reduced ($-$1) in individuals whose T2D genetic architecture is dominated by each functional cluster.
A concordance score of $+1$ was assigned when the observed and expected directions agreed; $-1$ when they were opposite (discordant); and $0$ when no directional expectation could be established (neutral).
Genomic loci were always assigned concordance $+1$ because the pipeline retains only OR\,$>$\,1 variants (risk-elevating direction), which is by construction concordant with the cluster's genetic risk profile.

Example discordant assignments that motivated this filter: in the cardiovascular-enriched cohort (\texttt{cardio\_product}), waist circumference, hip circumference, and basal metabolic rate were all negatively associated with T2D case status ($r < 0$), yet their genome-wide loci map into the obesity cluster.
This pattern---reduced adiposity phenotypes alongside genetic loading at obesity-related loci---is consistent with a lean insulin-resistant phenotype (mechanistically proximate to the lipodystrophy and body-fat clusters of \citet{Suzuki2024}) rather than classical obesity, and is analogous to the ``fat-distribution without obesity'' phenotype described by \citet{Yaghootkar2014}.

Concordance scores were propagated to the subtype-assignment step: for each functional cluster within a cohort, the net concordance score per putative T2D subtype was computed as the sum of feature-level concordance values (concordant~$= 1.0$, neutral~$= 0.5$, discordant~$= 0.0$).
Only subtypes with a positive net concordance score were retained in the circos visualisation and the cohort-enrichment pie charts.
This ensures that subtype annotations reflect the direction of the underlying phenotypic signal rather than purely the genomic locus membership of each feature.

\subsection*{Assignment of Functional Clusters to T2D Subtypes}

Functional clusters were assigned to T2D subtypes using a curated evidence map (\texttt{cluster\_subtype\_map.csv}) drawing on published subtype characterisation studies.

The five-subtype taxonomy of \citet{Ahlqvist2018}---severe autoimmune diabetes (SAID), severe insulin-deficient diabetes (SIDD), severe insulin-resistant diabetes (SIRD), mild obesity-related diabetes (MOD), and mild age-related diabetes (MARD)---provided the primary classification framework.
In accordance with the literature, LADA was considered mechanistically equivalent to SAID (autoimmune beta-cell destruction) and MODY was considered an extreme of the SIDD spectrum (monogenic beta-cell dysfunction); neither was retained as a distinct subtype category.

Cluster-to-subtype assignments were based on:
\begin{itemize}
  \item \citet{Yaghootkar2014} and \citet{Kim2023}, who demonstrated that lipodystrophy-related loci (IRS1, GRB14, PPARG) and body fat distribution loci segregate with SIRD through fat-distribution-mediated insulin resistance;
  \item \citet{Ahlqvist2018} and \citet{Kim2023}, who showed that beta-cell clusters (beta-cell~1, beta-cell~2, proinsulin, and residual glycaemic) are enriched in SIDD and, for shared autoimmune loci, in SAID;
  \item \citet{Aly2023}, who characterised the metabolic syndrome and hyper-insulin secretion clusters as predominantly MOD through compensatory hyperinsulinaemia in the context of insulin resistance;
  \item \citet{Suzuki2024}, whose cluster clinical phenotype profiles (HbA\textsubscript{1c}, HOMA-2B, HOMA-2IR, BMI, triglycerides, eGFR) were used directly as quantitative evidence, with subtype association declared where the standardised $z$-score of the clinical mean exceeded a threshold of $z \geq 0.5$ relative to the cross-subtype distribution.
\end{itemize}

HLA-region features were handled separately: consistent with \citet{Mansour2023}, only alleles at \textit{HLA-DQB1} and \textit{HLA-DQA1} loci were assigned to SAID; all other HLA alleles (\textit{HLA-A}, \textit{-B}, \textit{-C}, \textit{-DPA1}, \textit{-DPB1}, \textit{-DRB1}, \textit{-DRB4}, \textit{-DRB5}) were retained in the HLA-autoimmune tracking cluster without an automatic subtype assignment.

Clusters for which direct published evidence was unavailable (ALP-negative, SHBG/LpA, lipoprotein~A, cholesterol, bilirubin) were flagged for manual review (\texttt{assign\_manually = 1}) and excluded from automated subtype enrichment calculations.

Direct evidence from Suzuki~\textit{et al.}\ clinical phenotype tables and Mansour~\textit{et al.}\ HLA locus data were incorporated as additional feature-level subtype assignments via files matching the pattern \texttt{subtype*\_map.csv} in the subtype research directory, providing 27 and 2 feature--subtype records, respectively.
These direct assignments supplemented but did not override cluster-level assignments.

The relative contribution of each PRS cohort to each T2D subtype was summarised using concordance-weighted feature counts normalised to the subtype total, and visualised as donut pie charts (one per subtype).
Cohort contributions that were entirely discordant (net concordance score $\leq 0$) were excluded from these charts, so each wedge size reflects only the evidence supporting—rather than contradicting—a given subtype's phenotypic profile within that cohort.
